\section{Gerbang Universal dan Kelengkapan Fungsional}

Konsep kelengkapan fungsional (\textit{functional completeness}) merupakan salah satu hasil sentral dalam teori rangkaian digital. Suatu himpunan gerbang logika dikatakan lengkap secara fungsional jika setiap fungsi Boolean dapat diimplementasikan hanya menggunakan gerbang-gerbang dalam himpunan tersebut. Secara praktis, gerbang NAND secara tunggal bersifat lengkap secara fungsional, dan demikian pula gerbang NOR. Kedua gerbang ini disebut \logic{gerbang universal} \cite{rosen2019}.

\subsection{Pembuktian Universalitas NAND}

Untuk membuktikan bahwa NAND bersifat universal, cukup ditunjukkan bahwa ketiga operasi dasar (NOT, AND, OR) dapat diimplementasikan hanya menggunakan gerbang NAND. Jika $\uparrow$ melambangkan operasi NAND, maka $(A \uparrow B) = (A \cdot B)'$.

\begin{enumerate}[label=\textbf{\arabic*.}]
  \item \textbf{NOT dari NAND}: Hubungkan kedua masukan NAND ke sinyal yang sama:
  \[ A \uparrow A = (A \cdot A)' = A' \]

  \item \textbf{AND dari NAND}: AND adalah komplemen dari NAND, sehingga dibutuhkan dua gerbang NAND:
  \[ A \cdot B = ((A \cdot B)')' = (A \uparrow B) \uparrow (A \uparrow B) \]

  \item \textbf{OR dari NAND}: Berdasarkan hukum De Morgan, $A + B = (A' \cdot B')'$:
  \[ A + B = (A' \cdot B')' = (A \uparrow A) \uparrow (B \uparrow B) \]
  Diperlukan tiga gerbang NAND: dua untuk inversi dan satu untuk operasi NAND pada hasil inversi.
\end{enumerate}

\begin{figure}[H]
\centering
\begin{tikzpicture}[circuit logic US, huge circuit symbols,
                    every circuit symbol/.style={fill=cyan!10, draw=cyan!50!black, thick},
                    scale=0.85, transform shape]

  % NOT from NAND
  \node[font=\bfseries] at (0, 1.5) {NOT dari NAND};
  \node[nand gate] (n1) at (1.5, 0) {};
  \draw (n1.input 1) -- ++(-0.5,0) coordinate (j1);
  \draw (n1.input 2) -- ++(-0.5,0) coordinate (j2);
  \draw (j1) -- ++(-0.5, 0) |- (j2);
  \draw (j1) -- ++(-0.5,0) -- ++(-0.3,0) node[left] {$A$};
  \draw (n1.output) -- ++(0.5,0) node[right] {$A'$};

  % AND from NAND
  \node[font=\bfseries] at (6, 1.5) {AND dari NAND};
  \node[nand gate] (n2) at (7, 0) {};
  \node[nand gate] (n3) at (9.5, 0) {};
  \draw (n2.input 1) -- ++(-0.5,0) node[left] {$A$};
  \draw (n2.input 2) -- ++(-0.5,0) node[left] {$B$};
  \draw (n2.output) -- (n3.input 1);
  \draw (n2.output) ++(0.3,0) |- (n3.input 2);
  \draw (n3.output) -- ++(0.5,0) node[right] {$AB$};

\end{tikzpicture}
\caption{Implementasi NOT dan AND menggunakan Gerbang NAND}
\label{fig:nand-universal}
\end{figure}

\subsection{Pembuktian Universalitas NOR}

Dengan cara yang analog, universalitas NOR dibuktikan dengan mengimplementasikan NOT, OR, dan AND hanya menggunakan gerbang NOR. Jika $\downarrow$ melambangkan operasi NOR, maka $(A \downarrow B) = (A + B)'$.

\begin{enumerate}[label=\textbf{\arabic*.}]
  \item \textbf{NOT dari NOR}: $A \downarrow A = (A + A)' = A'$
  \item \textbf{OR dari NOR}: $A + B = ((A + B)')' = (A \downarrow B) \downarrow (A \downarrow B)$
  \item \textbf{AND dari NOR}: $A \cdot B = (A' + B')' = (A \downarrow A) \downarrow (B \downarrow B)$
\end{enumerate}

\subsection{Perbandingan Gerbang yang Bersifat Universal}

Tidak semua gerbang bersifat universal. Tabel~\ref{tab:universalitas} menunjukkan bahwa dari tujuh jenis gerbang logika, hanya NAND dan NOR yang secara individual bersifat lengkap secara fungsional.

\begin{table}[H]
\centering
\renewcommand{\arraystretch}{1.3}
\begin{tabular}{|l|c|l|}
\hline
\textbf{Gerbang} & \textbf{Universal?} & \textbf{Alasan} \\ \hline
AND & Tidak & Tidak dapat menghasilkan NOT \\ \hline
OR & Tidak & Tidak dapat menghasilkan NOT \\ \hline
NOT & Tidak & Tidak dapat menghasilkan AND/OR (hanya uner) \\ \hline
\textbf{NAND} & \textbf{Ya} & Dapat menghasilkan NOT, AND, OR \\ \hline
\textbf{NOR} & \textbf{Ya} & Dapat menghasilkan NOT, AND, OR \\ \hline
XOR & Tidak & Tidak dapat menghasilkan AND \\ \hline
XNOR & Tidak & Tidak dapat menghasilkan AND \\ \hline
\end{tabular}
\caption{Status Kelengkapan Fungsional Gerbang Logika}
\label{tab:universalitas}
\end{table}

\subsection{Implikasi Praktis dalam Produksi}

Universalitas NAND dan NOR memiliki konsekuensi praktis yang besar bagi industri semikonduktor. Dalam proses fabrikasi chip (\textit{wafer fabrication}), memproduksi satu jenis gerbang secara massal jauh lebih efisien dibandingkan memproduksi beragam jenis gerbang pada satu chip. Keluarga logika TTL (\textit{Transistor-Transistor Logic}) menggunakan NAND sebagai komponen primitif, sementara keluarga logika ECL (\textit{Emitter-Coupled Logic}) menggunakan NOR \cite{electronics_tutorials_gates}.

Selain efisiensi fabrikasi, penggunaan satu jenis gerbang juga menyederhanakan proses verifikasi dan pengujian. Alat-alat EDA modern secara otomatis mengonversi desain dari bentuk AND-OR-NOT ke dalam bentuk NAND-only atau NOR-only melalui proses yang disebut \logic{technology mapping}. Proses ini merupakan bagian integral dari alur desain VLSI (\textit{Very Large Scale Integration}) yang menghasilkan chip-chip terintegrasi modern.
